% arXiv-compliant paper skeleton for sia_rust
% Based on kourgeorge/arxiv-style (NeurIPS-inspired preprint style)
% Compile with: pdflatex main && pdflatex main && pdflatex main

\documentclass{article}
\usepackage{arxiv}

\usepackage[utf8]{inputenc}
\usepackage[T1]{fontenc}
\usepackage{hyperref}
\usepackage{url}
\usepackage{booktabs}
\usepackage{amsmath,amssymb,amsfonts}
\usepackage{graphicx}
\usepackage{algorithm}
\usepackage{algpseudocode}

\title{Self-Improving AI Agents in Rust: A Native Implementation of Harness and Weight Updates with Adaptive Scheduling}

\author{
  Micah Stubbs \thanks{Corresponding author} \\
  \texttt{micah.stubbs@gmail.com} \\
  \small Independent Researcher / CaseMirror.ai
}

\date{\today}

\begin{document}

\maketitle

\begin{abstract}
We present sia\_rust, a high-performance, memory-safe reimplementation of the Self-Improving AI (SIA) framework in Rust. Building on the recent SIA paper (Hebbar et al., arXiv:2605.27276), we unify harness (scaffold) evolution with test-time weight adaptation in a single Feedback-Agent loop. Our key contributions include: (1) a native Rust LLM agent runner built on rig-core with full trajectory logging middleware, (2) an adaptive harness/weight scheduler that learns when to apply each improvement lever, (3) ``SIA Studio'', a real-time interactive visualization dashboard for live self-improvement demos, and (4) a production-grade security model and threat analysis for self-modifying agent systems. We demonstrate measurable gains on legal reasoning, GPU kernel optimization, and scientific computing tasks while maintaining strict safety invariants. The system is designed for Nebius H200-scale deployments and is positioned as both a practical framework and a research platform for recursive self-improvement. All code, trajectories, and evaluation harnesses are open source.
\end{abstract}

\keywords{Self-improving agents, Recursive self-improvement, LLM agents, Rust, rig-core, Meta-RL, Trajectory logging, Verifier robustness}

\section{Introduction}

The recent release of the SIA paper (Hebbar et al., arXiv:2605.27276) demonstrated that unifying harness evolution with weight updates in a single Feedback-Agent loop yields significant gains across legal, systems, and scientific domains. However, the original implementation was Python-centric, limiting performance, safety guarantees, and ease of deployment on modern GPU infrastructure.

In this work we present \textbf{sia\_rust}, a complete native Rust port and extension of the SIA framework. Our implementation prioritizes:
\begin{itemize}
    \item Memory safety and deterministic execution for self-modifying code
    \item High-fidelity trajectory capture as the substrate for intelligent improvement decisions
    \item First-class support for sponsor infrastructure (Nebius H200 GPUs)
    \item A live, judge-friendly demonstration experience (``SIA Studio'')
    \item Clear path toward the paper's own stated future direction: meta-RL over the harness-vs-weight decision policy
\end{itemize}

We make the following contributions:
\begin{enumerate}
    \item A production-ready native Rust LLM client abstraction on top of rig-core with Anthropic, OpenAI, and other providers.
    \item Trajectory logging middleware that captures full execution traces (prompts, tool calls, errors, verifier signals) for downstream analysis.
    \item An adaptive scheduler (heuristic + multi-armed bandit foundation) that learns when to trigger harness updates versus weight updates.
    \item ``SIA Studio'' --- an interactive Axum-based web dashboard for real-time visualization of self-improvement.
    \item Security hardening and a formal threat model tailored to self-modifying agent systems.
    \item Robust verifier contracts and domain task templates suitable for Applied and Research tracks.
\end{enumerate}

\section{Related Work}

[To be expanded with citations from the original SIA paper, recent self-improving agent literature, rig-core ecosystem, and meta-RL / bandit work on LLM improvement levers.]

\section{System Architecture}

[Describe the three-agent loop (Meta \to Target \to Feedback), the Python bridge (transitional), native rig-core path, trajectory capture points, and how the adaptive scheduler sits in the Feedback-Agent. Reference issues #62, #51, #46, #65.]

\section{Adaptive Harness/Weight Scheduler}

[Detail the heuristic baseline, UCB/Thompson Sampling bandit layer, contextual features from trajectory diagnosis (Chen et al., arXiv:2606.06324), and logging for future meta-RL. This directly implements a future direction called out in Hebbar et al. arXiv:2605.27276 p.13.]

\section{Experimental Results}

[Report gains on LawBench subset, GPU kernel optimization, scRNA-seq style tasks, and any new hackathon tasks. Include ablation on scheduler policy, trajectory fidelity impact, and Nebius H200 scaling numbers.]

\section{Security and Threat Model}

[Document the capability model, sandboxing strategy (Docker + future WASI/landlock), and formal analysis of risks unique to self-modifying systems.]

\section{Demonstration: SIA Studio}

[Describe the live dashboard, real-time charts, trajectory replay, one-click generation runs, and how it enables compelling 3--4 minute hackathon demos.]

\section{Conclusion and Future Work}

[Reiterate contributions, position sia\_rust as both a winning hackathon artifact and a serious research platform, and outline the path to full meta-RL over the improvement policy.]

\section*{Acknowledgments}

We thank the organizers of the Self-Improving AI Agents hackathon (AGI House SF + Nebius), the authors of the original SIA paper, and the rig-core and broader Rust LLM ecosystem contributors.

\bibliographystyle{plain}
\begin{thebibliography}{99}

\bibitem{hebbar2026sia}
Hebbar et al. ``SIA: Self Improving AI with Harness \& Weight Updates''. arXiv:2605.27276, 2026.

% Add more references as sections are filled (Chen et al. 2026, Liu et al. 2025, rig-core papers, etc.)

\end{thebibliography}

\end{document}